\chapter{Fazit und Ausblick}
\textbf{Johanna Hoppe}

Ziel dieser Arbeit war es, die synaptische Lokalisation von Caspr2 sowie den Einfluss
von aAK gegen Caspr2 auf dessen Interaktion mit Contactin-2 zu
untersuchen. Dabei standen zwei Fragestellungen im Mittelpunkt. Zum einen sollte
herausgefunden werden, ob Caspr2 bevorzugt an exzitatorischen oder inhibitorischen
Synapsen lokalisiert ist. Zum anderen wurde untersucht, ob AK gegen die
Discoidin- bzw. Laminin-G1-Domäne von Caspr2 dessen Kolokalisation mit
Contactin-2 beeinflussen.

Zur Untersuchung der synaptischen Lokalisation wurde die Kolokalisation von Caspr2
mit den synaptischen Markern VGat und VGlut betrachtet. Dabei konnte kein
statistisch signifikanter Unterschied zwischen den beiden Markern festgestellt werden.
Eine eindeutige Aussage über eine bevorzugte Lokalisation von Caspr2 an
exzitatorischen oder inhibitorischen Synapsen ist anhand der vorliegenden Ergebnisse
nicht möglich. Insbesondere die fehlende spezifische Bindung der verwendeten
Caspr2-Primär-AK sowie die begrenzte Anzahl auswertbarer Messwerte
schränken die Aussagekraft der Untersuchung ein.

Auch bei der Untersuchung des Einflusses der AK gegen die Discoidin- und
Laminin-G1-Domäne konnte kein statistisch signifikanter Unterschied zwischen den
Gruppen nachgewiesen werden. Da aufgrund einer fehlgeschlagenen Kontrollfärbung
keine unbehandelte Vergleichsgruppe vorlag, lässt sich nicht beurteilen, ob die
AK-Behandlung die Kolokalisation von Caspr2 und Contactin-2 tatsächlich
verändert.

Zusammenfassend konnten die durchgeführten Experimente keine abschließenden
Antworten auf die beiden Fragestellungen liefern. Die Ergebnisse zeigen jedoch, dass
die Aussagekraft der Untersuchungen maßgeblich von der Spezifität der verwendeten
AK, der Qualität der Färbungen und einer ausreichenden Datengrundlage
abhängt. 

Um die Fragestellungen beantworten zu können, sollten zunächst die verwendeten
AK auf ihre Spezifität überprüft werden. Dies ist durch erneute Färbungen von Caspr2 mit anti-Caspr2-aAK durch Dr. Josefine Sell bereits erfolgt (siehe Kapitel 8 Abbildung \ref{fig:selltest}).
Zur vollständigen Auswertung des zweiten Experiments muss die misslungene Kontrollprobe aufgrund des Fehlers im Protokoll angepasst wiederholt werden. Es muss also, wie in Kapitel 8 beschrieben, eine zusätzliche Gegenfärbung von Caspr2 durchgeführt werden, um eine veränderte Kolokalisation zu Contactin-2 untersuchen zu können. Mit einer größeren Anzahl an Messungen könnten zudem belastbarere Aussagen zu den Fragestellungen getroffen werden. 
%Auf dieser Grundlage könnte erneut untersucht werden, ob Caspr2 eine bevorzugte synaptische Lokalisation aufweist und ob Autoantikörper gegen unterschiedliche Domänen von Caspr2 dessen räumliche Beziehung zu Contactin-2 beeinflussen.